\definicion{Modelo cliente-servidor}{Arquitectura web donde el cliente (navegador) solicita recursos y el servidor los aloja y responde, complementando lenguajes de cliente (HTML, CSS, JS) y de servidor.}
\subtituloA{Componentes del navegador}

\lista{
  \textbf{Arquitectura de referencia:} El navegador incluye ocho componentes, como interfaz, buscador, renderizado, comunicaciones, intérprete JS, parser XML, visualización y persistencia. |
  \textbf{Integración de scripts:} Se usa la etiqueta <script> con rutas relativas, absolutas o remotas, complementada con <noscript> y sintaxis CDATA en XHTML.
}
\nota{Para el examen}{
Recordar que JavaScript es la capa de comportamiento de la web y que el sitio debe funcionar sin depender de él. El flujo es: editor $\rightarrow$ prueba en navegadores $\rightarrow$ validación W3C.
}
